\documentclass[11pt,a4paper]{article}

\usepackage[T1]{fontenc}
\usepackage[utf8]{inputenc}
\usepackage{lmodern}
\usepackage{amssymb,amsmath,amsthm,mathtools}
\usepackage[margin=2.5cm]{geometry}
\usepackage{parskip}
\usepackage[hidelinks]{hyperref}
\usepackage{microtype}

% Theorem environments (arXiv / journal style)
\theoremstyle{plain}
\newtheorem{theorem}{Theorem}[section]
\newtheorem{proposition}[theorem]{Proposition}
\newtheorem{lemma}[theorem]{Lemma}
\newtheorem{corollary}[theorem]{Corollary}

\theoremstyle{definition}
\newtheorem{definition}[theorem]{Definition}
\newtheorem{example}[theorem]{Example}
\newtheorem{remark}[theorem]{Remark}

% Notation
\newcommand{\Ext}{\mathrm{Ext}}
\newcommand{\Shape}{\mathrm{Shape}}
\newcommand{\Pos}{\mathrm{Pos}}
\newcommand{\Type}{\mathbf{Type}}
\newcommand{\id}{\mathrm{id}}
\newcommand{\Ihom}{\mathrm{Ihom}}
\newcommand{\Cont}{\mathbf{Cont}}
\newcommand{\Hom}{\mathrm{Hom}}

\title{Closure of the Dirichlet Tensor on Containers}
\author{Extracted from a machine-checked development}
\date{}

\begin{document}
\maketitle

\begin{abstract}
We prove that the Dirichlet tensor $\otimes$ on the category $\Cont$ of
containers is \emph{closed}: for every container $q$ the functor
$- \otimes q$ has a right adjoint $\Ihom(q,-)$. The internal hom is itself
a container, whose shapes are the container morphisms $q \Rightarrow r$ and
whose positions over a morphism $(u,g)$ are pairs consisting of a $q$-shape
and an $r$-position over its image. The adjunction is exhibited by an
explicit currying bijection of hom-sets, with both round trips and both
naturalities holding by pure rearrangement of dependent products. The proofs
of all stated results have been machine-checked; we record only the
mathematical content.
\end{abstract}

\tableofcontents
\bigskip

%======================================================================
\section{Introduction}
%======================================================================

Containers form a category $\Cont$ whose morphisms are pairs of a forward
map on shapes and a backward map on positions. Among the monoidal structures
on $\Cont$ is the \emph{Dirichlet tensor}: for containers $C = (S,P)$ and
$D = (T,Q)$,
\begin{equation}
  C \otimes D \;=\; \bigl(S \times T,\; (s,t) \mapsto P(s) \times Q(t)\bigr).
\end{equation}
Shapes are paired; positions are products. This is dual in spirit to the
categorical product of containers, which uses the same shape product but
coproducts of positions.

A monoidal category is closed when each functor $- \otimes q$ admits a
right adjoint. In the polynomial-functor presentation this is classical
(Niu--Spivak); the present note records the corresponding statement in the
container presentation, with fully explicit currying and uncurrying maps
and with naturality in both variables.

The material is elementary and self-contained. We work throughout in a
single universe of types, writing $\Type$ for that universe. Morphisms of
containers and the Dirichlet tensor are assumed as in the prior notes on
containers and monoidal structure; we recall only what is needed. Proofs
are indicated by the phrase ``by lean,'' meaning they have been checked by
a proof assistant rather than written out in prose.

\paragraph{What this note claims.}
In ordinary mathematical language, the claim is: the Dirichlet monoidal
structure $(\Cont, \otimes, y)$ is closed. The internal hom
$\Ihom(q,r)$ is the container whose shape type is literally the set of
container morphisms $q \Rightarrow r$, and the currying bijection
\begin{equation}
  \Hom(p \otimes q,\, r)
  \;\cong\;
  \Hom\bigl(p,\, \Ihom(q,r)\bigr)
\end{equation}
is pure currying of dependent products---split across the two components of
the backward position map of a morphism $p \otimes q \Rightarrow r$. There
is no extra data, no transport, and no exotic naturality constraint beyond
what the container presentation already forces.

%======================================================================
\section{Recollections}
%======================================================================

\begin{definition}[Container and morphism]
A \textbf{container} $C$ consists of a type $\Shape(C)$ of shapes and, for
each shape $s$, a type $\Pos(C)(s)$ of positions. A \textbf{morphism}
$\varphi : C \Rightarrow D$ consists of a shape map
$u : \Shape(C) \to \Shape(D)$ and, for each $s$, a position map
$g_s : \Pos(D)(u(s)) \to \Pos(C)(s)$ (contravariant on positions). We write
$\Hom(C,D)$ for the type of such morphisms.
\end{definition}

\begin{definition}[Dirichlet tensor]\label{def:dirichlet}
Let $C = (S,P)$ and $D = (T,Q)$ be containers. Their \textbf{Dirichlet
tensor} is the container
\begin{equation}
  C \otimes D \;=\;
  \bigl(
    S \times T,\;
    (s,t) \mapsto P(s) \times Q(t)
  \bigr).
\end{equation}
The unit $y$ is the one-shape, one-position container (the same object that
serves as unit for sequential composition of containers).
\end{definition}

\begin{definition}[Bifunctorial action]\label{def:dir2}
Given morphisms $\mu : C \Rightarrow C'$ and $\nu : D \Rightarrow D'$, the
\textbf{tensor of morphisms}
$\mu \otimes \nu : C \otimes D \Rightarrow C' \otimes D'$ acts
componentwise:
\begin{align}
  (\mu \otimes \nu)\text{ on shapes:}
  &\quad (s,t) \mapsto \bigl(u_\mu(s),\, u_\nu(t)\bigr), \\
  (\mu \otimes \nu)\text{ on positions:}
  &\quad (p,q) \mapsto \bigl(g_\mu(p),\, g_\nu(q)\bigr).
\end{align}
\end{definition}

%======================================================================
\section{The internal hom}
%======================================================================

The internal hom for $\otimes$ is a container whose shapes \emph{are}
morphisms. That is not a notational convenience: it is the content of the
closed structure (cf.\ Niu--Spivak, Eq.~(4.79), transported to containers).

\begin{definition}[Internal hom]\label{def:ihom}
Let $q$ and $r$ be containers. The \textbf{internal hom}
$\Ihom(q,r)$ is the container defined by
\begin{align}
  \Shape\bigl(\Ihom(q,r)\bigr)
  &\;=\;
  \Hom(q,r), \\[4pt]
  \Pos\bigl(\Ihom(q,r)\bigr)(f)
  &\;=\;
  \sum_{t \,:\, \Shape(q)}\;
  \Pos(r)\bigl(u_f(t)\bigr),
\end{align}
where $f = (u_f, g_f) : q \Rightarrow r$. Thus a shape of $\Ihom(q,r)$ is a
container morphism $q \Rightarrow r$, and a position over such a shape is a
pair $\langle t,\, i \rangle$ consisting of a $q$-shape $t$ and an
$r$-position $i$ over the image shape $u_f(t)$.
\end{definition}

\begin{proposition}[Shapes are morphisms]\label{prop:ihom-shape}
For all containers $q$ and $r$,
\begin{equation}
  \Shape\bigl(\Ihom(q,r)\bigr) \;=\; \Hom(q,r)
\end{equation}
definitionally.
\end{proposition}

\begin{proof}
By lean.
\end{proof}

\begin{proposition}[Positions of the internal hom]\label{prop:ihom-pos}
For every morphism $f : q \Rightarrow r$,
\begin{equation}
  \Pos\bigl(\Ihom(q,r)\bigr)(f)
  \;=\;
  \sum_{t \,:\, \Shape(q)}\; \Pos(r)\bigl(u_f(t)\bigr)
\end{equation}
definitionally.
\end{proposition}

\begin{proof}
By lean.
\end{proof}

\begin{remark}
That the shape type is the hom-set of $\Cont$ is what makes the adjunction
\emph{internal}: the internal hom is a container, not merely a set of
morphisms. The extra data---the position family---is precisely what is
needed for the currying bijection of the next section to be a bijection of
container morphisms rather than of bare functions.
\end{remark}

%======================================================================
\section{Currying and uncurrying}
%======================================================================

Unfolded, a morphism $F : p \otimes q \Rightarrow r$ consists of
\begin{itemize}
  \item a shape map $u_F : \Shape(p) \times \Shape(q) \to \Shape(r)$, and
  \item a position map
        $g_F : (s,t) \to \Pos(r)\bigl(u_F(s,t)\bigr)
        \to \Pos(p)(s) \times \Pos(q)(t)$.
\end{itemize}
A morphism $G : p \Rightarrow \Ihom(q,r)$ consists of
\begin{itemize}
  \item a shape map $u_G : \Shape(p) \to \Hom(q,r)$, sending each $s$ to a
        pair $(u_s, \varphi_s)$, and
  \item a position map
        $g_G : (s : \Shape(p)) \to
        \bigl(\sum_t \Pos(r)(u_s(t))\bigr) \to \Pos(p)(s)$.
\end{itemize}
The bijection is pure currying, \emph{split across the two components of
$g_F$}: the $\Pos(q)(t)$ component of $g_F$ becomes the backward position
map $\varphi_s$ of the shape $u_G(s)$, and the $\Pos(p)(s)$ component
becomes $g_G$. That split is the whole content---it is why the positions
of $\Ihom(q,r)$ are $\sum_t \Pos(r)(u(t))$ and not something else.

\begin{definition}[Currying]\label{def:curry}
Let $F : p \otimes q \Rightarrow r$ be a container morphism. The
\textbf{curried morphism}
\begin{equation}
  \mathrm{curry}(F) \;:\; p \;\Rightarrow\; \Ihom(q,r)
\end{equation}
is defined by
\begin{align}
  u_{\mathrm{curry}(F)}(s)
  &\;=\;
  \bigl(
    t \mapsto u_F(s,t),\;
    (t,i) \mapsto \pi_2\bigl(g_F(s,t)(i)\bigr)
  \bigr), \\[4pt]
  g_{\mathrm{curry}(F)}\bigl(s,\, \langle t, i \rangle\bigr)
  &\;=\;
  \pi_1\bigl(g_F(s,t)(i)\bigr).
\end{align}
Thus the $q$-component of $F$'s backward position map is packaged
\emph{into the shape} of the internal hom, while the $p$-component becomes
the backward position map of the curried morphism.
\end{definition}

\begin{definition}[Uncurrying]\label{def:uncurry}
Let $G : p \Rightarrow \Ihom(q,r)$ be a container morphism. The
\textbf{uncurried morphism}
\begin{equation}
  \mathrm{uncurry}(G) \;:\; p \otimes q \;\Rightarrow\; r
\end{equation}
is defined by
\begin{align}
  u_{\mathrm{uncurry}(G)}(s,t)
  &\;=\;
  u_{u_G(s)}(t), \\[4pt]
  g_{\mathrm{uncurry}(G)}\bigl((s,t),\, i\bigr)
  &\;=\;
  \Bigl(
    g_G\bigl(s,\, \langle t, i \rangle\bigr),\;
    g_{u_G(s)}(t,i)
  \Bigr).
\end{align}
Read the shape map of $G$ at $s$ as a morphism $q \Rightarrow r$; its shape
map supplies the $r$-shape, and its backward position map supplies the
$\Pos(q)(t)$ component. The $\Pos(p)(s)$ component is $G$'s own backward
position map applied to the position $\langle t, i \rangle$.
\end{definition}

\begin{theorem}[Uncurry after curry]\label{thm:uncurry-curry}
For every morphism $F : p \otimes q \Rightarrow r$,
\begin{equation}
  \mathrm{uncurry}\bigl(\mathrm{curry}(F)\bigr) \;=\; F.
\end{equation}
\end{theorem}

\begin{proof}
By lean.
\end{proof}

\begin{theorem}[Curry after uncurry]\label{thm:curry-uncurry}
For every morphism $G : p \Rightarrow \Ihom(q,r)$,
\begin{equation}
  \mathrm{curry}\bigl(\mathrm{uncurry}(G)\bigr) \;=\; G.
\end{equation}
\end{theorem}

\begin{proof}
By lean.
\end{proof}

\begin{remark}[Why both round trips are immediate]
Both identities are rearrangements of a $\Sigma$/$\times$ telescope. The only
ingredients are $\eta$ for products (rebuilding a pair from its projections),
$\eta$ for dependent sums (rebuilding $\langle t, i \rangle$), and $\eta$
for the structure of a container morphism (rebuilding a morphism from its
shape and position components). No transport of paths, no function
extensionality beyond definitional $\eta$, and no appeal to naturality are
required. The same definitional character holds for the monoidal coherences
of $\otimes$ itself: the Dirichlet tensor is well behaved at the level of
definitional equality in a way that sequential composition and the directed
container comonad laws are not.
\end{remark}

%======================================================================
\section{The closure theorem}
%======================================================================

\begin{definition}[Type equivalence]
A \textbf{type equivalence} $A \simeq B$ is a pair of maps
$f : A \to B$ and $g : B \to A$ together with proofs that $g \circ f = \id_A$
and $f \circ g = \id_B$.
\end{definition}

\begin{theorem}[Dirichlet closure]\label{thm:closure}
For all containers $p$, $q$, and $r$ there is a type equivalence
\begin{equation}
  \Hom(p \otimes q,\, r)
  \;\simeq\;
  \Hom\bigl(p,\, \Ihom(q,r)\bigr),
\end{equation}
with forward map $\mathrm{curry}$ and inverse $\mathrm{uncurry}$. The two
round-trip identities are Theorems~\ref{thm:uncurry-curry}
and~\ref{thm:curry-uncurry}. Equivalently,
\begin{equation}
  - \otimes q \;\dashv\; \Ihom(q,-).
\end{equation}
\end{theorem}

\begin{proof}
By lean.
\end{proof}

\begin{remark}
Together with the monoidal structure $(\Cont, \otimes, y)$---associator,
unitors, pentagon, and triangle---this makes $(\Cont, \otimes, y)$ a
\textbf{closed monoidal category}, with internal hom $\Ihom$.
\end{remark}

%======================================================================
\section{Functoriality of the internal hom and naturality}
%======================================================================

The currying bijection is natural in $p$ (contravariantly, by
precomposition) and in $r$ (covariantly, by postcomposition). Naturality in
$r$ requires the action of $\Ihom(q,-)$ on morphisms, which we record first.

\begin{definition}[Action on morphisms]\label{def:ihom-map}
Let $\psi : r \Rightarrow r'$ be a container morphism. The induced morphism
\begin{equation}
  \Ihom(q,\psi) \;:\; \Ihom(q,r) \;\Rightarrow\; \Ihom(q,r')
\end{equation}
is defined by
\begin{align}
  \text{on shapes:}
  &\quad f \mapsto f \circ \psi
  \qquad\text{(postcomposition)}, \\[4pt]
  \text{on positions:}
  &\quad
  \langle t,\, i \rangle
  \mapsto
  \bigl\langle t,\;
    g_\psi\bigl(u_f(t)\bigr)(i)
  \bigr\rangle.
\end{align}
A position $\langle t, i \rangle$ of $\Ihom(q,r')$ over the shape
$f \circ \psi$ has $i$ an $r'$-position over $\psi$'s image of $u_f(t)$;
pulling $i$ back through $g_\psi$ yields an $r$-position, hence a position
of $\Ihom(q,r)$.
\end{definition}

\begin{proposition}[Functoriality of $\Ihom(q,-)$]\label{prop:ihom-func}
For every container $q$ and every container $r$,
\begin{equation}
  \Ihom(q, \id_r) \;=\; \id_{\Ihom(q,r)}.
\end{equation}
For morphisms $\psi : r \Rightarrow r'$ and $\psi' : r' \Rightarrow r''$,
\begin{equation}
  \Ihom(q,\, \psi \circ \psi')
  \;=\;
  \Ihom(q,\psi) \circ \Ihom(q,\psi').
\end{equation}
\end{proposition}

\begin{proof}
By lean.
\end{proof}

\begin{theorem}[Naturality in $p$]\label{thm:nat-left}
Let $h : p' \Rightarrow p$ and $F : p \otimes q \Rightarrow r$. Then
\begin{equation}
  \mathrm{curry}\bigl((h \otimes \id_q) \circ F\bigr)
  \;=\;
  h \circ \mathrm{curry}(F).
\end{equation}
\end{theorem}

\begin{proof}
By lean.
\end{proof}

\begin{theorem}[Naturality in $r$]\label{thm:nat-right}
Let $F : p \otimes q \Rightarrow r$ and $\psi : r \Rightarrow r'$. Then
\begin{equation}
  \mathrm{curry}(F \circ \psi)
  \;=\;
  \mathrm{curry}(F) \circ \Ihom(q,\psi).
\end{equation}
\end{theorem}

\begin{proof}
By lean.
\end{proof}

\begin{proposition}[Naturality of uncurry in $p$]\label{prop:uncurry-nat}
Let $h : p' \Rightarrow p$ and $G : p \Rightarrow \Ihom(q,r)$. Then
\begin{equation}
  \mathrm{uncurry}(h \circ G)
  \;=\;
  (h \otimes \id_q) \circ \mathrm{uncurry}(G).
\end{equation}
\end{proposition}

\begin{proof}
By lean.
\end{proof}

%======================================================================
\section{The counit: evaluation}
%======================================================================

Every adjunction has a counit. For Dirichlet closure the counit is
evaluation: applying a morphism stored as a shape of the internal hom.

\begin{definition}[Evaluation]\label{def:eval}
The \textbf{evaluation morphism}
\begin{equation}
  \mathrm{ev}_{q,r}
  \;:\;
  \Ihom(q,r) \otimes q \;\Rightarrow\; r
\end{equation}
is obtained by uncurrying the identity on $\Ihom(q,r)$:
\begin{equation}
  \mathrm{ev}_{q,r}
  \;=\;
  \mathrm{uncurry}\bigl(\id_{\Ihom(q,r)}\bigr).
\end{equation}
Unfolded, it sends the shape $(f,t)$ to $u_f(t)$---that is, it
\emph{applies} the morphism stored in the shape---and its backward position
map splits an $r$-position $i$ into the $\Ihom(q,r)$-position
$\langle t, i \rangle$ and the $q$-position $g_f(t,i)$.
\end{definition}

\begin{theorem}[Triangle identity / universal property]\label{thm:triangle}
For every morphism $F : p \otimes q \Rightarrow r$,
\begin{equation}
  \bigl(\mathrm{curry}(F) \otimes \id_q\bigr) \circ \mathrm{ev}_{q,r}
  \;=\;
  F.
\end{equation}
\end{theorem}

\begin{proof}
By lean.
\end{proof}

\begin{theorem}[Uniqueness of the curried mate]\label{thm:unique}
Let $F : p \otimes q \Rightarrow r$ and
$G : p \Rightarrow \Ihom(q,r)$. If
\begin{equation}
  (G \otimes \id_q) \circ \mathrm{ev}_{q,r} \;=\; F,
\end{equation}
then $G = \mathrm{curry}(F)$.
\end{theorem}

\begin{proof}
By lean.
\end{proof}

\begin{remark}
Existence (Theorem~\ref{thm:triangle}) is the statement that
$\mathrm{curry}(F)$ is a mate for $F$ under evaluation. Uniqueness
(Theorem~\ref{thm:unique}) is the observation that any $G$ satisfying the
triangle is recovered as $\mathrm{curry}$ of its composite with evaluation,
which is definitionally $\mathrm{uncurry}(G)$, so the claim reduces to the
round trip of Theorem~\ref{thm:curry-uncurry}.
\end{remark}

%======================================================================
\section{Discussion}
%======================================================================

\paragraph{What has been established.}
The Dirichlet tensor on containers is closed. The internal hom has shapes
equal to container morphisms and positions equal to pairs of a source shape
with a target position over its image. Currying is a pure split of the
backward position map of a morphism out of a tensor. Both round trips, both
naturalities, functoriality of $\Ihom(q,-)$, and the evaluation triangle
hold by definitional rearrangement.

\paragraph{Relation to the polynomial presentation.}
In the polynomial-functor (``Poly'') presentation of Niu--Spivak, the same
adjunction is Eq.~(4.79) of \emph{Polynomial Functors}. The container
presentation makes the split of position data completely explicit: one sees
exactly which component of the tensor's position product becomes shape data
of the internal hom and which becomes its position data.

\paragraph{Definitional character.}
Every proof in the development is by reflexivity of equality after unfolding
definitions. This is a feature of the Dirichlet tensor and of container
morphisms: both are built from products and dependent sums in a way that
definitional $\eta$ identifies the relevant telescopes. By contrast,
sequential composition of containers and the directed-container comonad laws
require genuine transports.

\paragraph{Scope.}
The development treats a single universe of types and does not address size
issues or universe polymorphism. It assumes the monoidal structure of
$\otimes$ (associator, unitors, coherences) from the companion note on the
Dirichlet tensor; those coherences are likewise definitional.

\paragraph{Machine checking.}
Every definition and theorem statement above corresponds to a declaration in
a machine-checked formalisation. The phrase ``by lean'' indicates that the
proof has been verified by a proof assistant; we have not reproduced the
formal proof scripts here. A mathematician reading this note need only
verify that the English statements match the intended mathematics; the
formal system is responsible for the derivations.

%======================================================================
\section*{Acknowledgements}
%======================================================================

The mathematical content follows the theory of polynomial functors and their
Dirichlet monoidal structure as developed by Niu and Spivak, transported to
the container presentation of Abbott, Altenkirch, and Ghani. The presentation
here is a prose rendering of a formal development intended for human review
of the theorem statements.

\begin{thebibliography}{9}

\bibitem{abbott2005}
M.~Abbott, T.~Altenkirch, and N.~Ghani.
\newblock Containers: constructing strictly positive types.
\newblock \emph{Theoretical Computer Science}, 342(1):3--27, 2005.

\bibitem{niu-spivak}
N.~Niu and D.~I.~Spivak.
\newblock Polynomial functors: a mathematical theory of interaction.
\newblock arXiv:2312.00990, 2023.

\end{thebibliography}

\end{document}
